\documentclass[12pt]{article}
\usepackage{amsmath,amssymb,amsthm}
\usepackage[margin=1in]{geometry}

\newtheorem{theorem}{Theorem}[section]
\newtheorem{lemma}[theorem]{Lemma}
\theoremstyle{definition}
\newtheorem{definition}[theorem]{Definition}
\newtheorem{exercise}[theorem]{Exercise}

\title{Principles of Mathematical Analysis: Chapter 5 Exercises}
\date{}

\begin{document}
\maketitle
\noindent The solutions below present the essential arguments in a concise, self-contained form.

Let $f$ be defined for all real $x$, and suppose that
\begin{align*}
\left\vert f(x)-f(y)\right\vert\le(x-y)^2
\end{align*}
for all $x$ and $y$. Prove that $f$ is constant.

\smallskip
\noindent\textbf{Solution.}\quad

Starting at $f(0)$, we show that $f(x)=f(0)$ for all $x\in \mathbb{R}$. Consider the case $x>0$, and for $n=1,2\ldots$, let $\{x_0,x_1,x_2,\ldots,x_n\}$ be a \textit{partition} of the interval $[0,x]$ with $x_0=0$, $x_n=x$, and $x_{j+1}-x_j=\frac{x}{n}$ for $j\in\{0,1,2,\ldots,n-1\}$. Then
\begin{align*}
\left\vert f(x)-f(0)\right\vert
&=\left\vert \sum_{j=0}^{n-1}f(x_{j+1})-f(x_j)\right\vert \\
&\le\sum_{j=0}^{n-1}\left\vert f(x_{j+1})-f(x_j)\right\vert \\
&\le\sum_{j=0}^{n-1}\left(x_{j+1}-x_j\right)^2 \\
&=\sum_{j=0}^{n-1}\left(\frac{x}{n}\right)^2 \\
&=\frac{x^2}{n}
\end{align*}
where the third line of the above inequality follows by assumption. Since $\frac{x^2}{n}\to 0$ as $n\to\infty$, we conclude that $\left\vert f(x)-f(0)\right\vert=0$ or $f(x)=f(0)$. The other case $x<0$ is in a similar way. Hence $f$ is constant.



\bigskip\noindent\hrule\bigskip

Suppose $f'(x)>0$ in $(a,b)$. Prove that $f$ is strictly increasing in $(a,b)$, and let $g$ be its inverse function. Prove that $g$ is differentiable, and that
\begin{align*}
g'(f(x))=\frac{1}{f'(x)}\quad(a<x<b)
\end{align*}

\smallskip
\noindent\textbf{Solution.}\quad

(i) By the mean value theorem (Theorem 5.10), for $x,y\in(a,b)$ with $x<y$,
\begin{align*}
f(y)-f(x)=f'(c)(y-x)>0
\end{align*}
for some $c\in(x,y)$. So $f$ is strictly increasing in $(a,b)$.
(ii) Let $g:f(a,b)\to(a,b)$ be the inverse function of $f$, i.e., $g(f(x))=x$ for all $x\in(a,b)$. We now show that
\begin{align*}
g'(y)=\lim_{z\to y}\frac{g(z)-g(y)}{z-y}
\end{align*}
exists for all $y\in f(a,b)$. Put $y=f(x)$ and $z=f(t)$, where $x,t\in(a,b)$, then since $f$ is continuous (by Theorem 5.2), so is $g$ (by Theorem 4.17), and $z\to y$ implies $t\to x$. It follows that
\begin{align*}
\lim_{z\to y}\frac{g(z)-g(y)}{z-y}
&=\lim_{t\to x}\frac{g(f(t))-g(f(x))}{f(t)-f(x)} \\
&=\lim_{t\to x}\frac{t-x}{f(t)-f(x)} \\
&=\lim_{t\to x}\frac{1}{\,\dfrac{f(t)-f(x)}{t-x}\,} \\
&=\dfrac{1}{\displaystyle\lim_{t\to x}\dfrac{f(t)-f(x)}{t-x}} \\
&=\frac{1}{f'(x)}
\end{align*}
Since $f'>0$, we see that $g'(y)$ is defined for all $y\in f(a,b)$. Hence $g$ is differentiable. It is also clear that $g'(f(x))=\frac{1}{f'(x)}$ for all $x\in(a,b)$.



\bigskip\noindent\hrule\bigskip

Suppose $g$ is a real function on $\mathbb{R}$, with bounded derivative (say $\left\vert g'\right\vert\le M$). Fix $\varepsilon>0$, and define $f(x)=x+\varepsilon g(x)$. Prove that $f$ is one-to-one if $\varepsilon$ is small enough. (A set of admissible value of $\varepsilon$ can be determined which depends only on $M$.)

\smallskip
\noindent\textbf{Solution.}\quad

If $0<\varepsilon<\frac{1}{M}$, then $g'\ge-M$ implies
\begin{align*}
f'(x)=1+\varepsilon g'(x)\ge 1-\varepsilon M>0
\end{align*}
Thus by Exercise 2, $f$ is strictly increasing, and hence $f$ is one-to-one.



\bigskip\noindent\hrule\bigskip

If
\begin{align*}
C_0+\frac{C_1}{2}+\cdots+\frac{C_{n-1}}{n}+\frac{C_n}{n+1}=0
\end{align*}
where $C_0,\ldots,C_n$ are real constants, prove that the equation
\begin{align*}
C_0+C_1x+\cdots+C_{n-1}x^{n-1}+C_nx^n=0
\end{align*}
has at least one real root between $0$ and $1$.

\smallskip
\noindent\textbf{Solution.}\quad

Let $f(x)=C_0x+\frac{C_1}{2}x^2+\cdots+\frac{C_n}{n+1}x^{n+1}$, then clearly $f$ is a differentiable real function (since $C_0,C_1,\ldots,C_n$ are real constants). Moreover, we have $f(0)=0$ and
\begin{align*}
f(1)=C_0+\frac{C_1}{2}+\cdots+\frac{C_{n-1}}{n}+\frac{C_n}{n+1}=0
\end{align*}
By the mean value theorem, there exists a point $x_0\in(0,1)$ such that $f'(x_0)=\frac{f(1)-f(0)}{1-0}=0$, i.e.,
\begin{align*}
C_0+C_1x_0+\cdots+C_{n-1}x_0^{n-1}+C_nx_0^n=0
\end{align*}



\bigskip\noindent\hrule\bigskip

Suppose $f$ is defined and differentiable for every $x>0$, and $f'(x)\to 0$ as $x\to+\infty$. Put $g(x)=f(x+1)-f(x)$. Prove that $g(x)\to 0$ as $x\to+\infty$.

\smallskip
\noindent\textbf{Solution.}\quad

By the mean value theorem, for $x>0$, let $c(x)$ be a real number such that $c(x)\in(x,x+1)$ and $f'(c(x))=f(x+1)-f(x)$. Then $c(x)\to+\infty$ as $x\to+\infty$, which implies
\begin{align*}
g(x)=f'(c(x))\to 0
\end{align*}



\bigskip\noindent\hrule\bigskip

Suppose
(a) $f$ is continuous for $x\ge 0$,
(b) $f'(x)$ exists for $x>0$,
(c) $f(0)=0$,
(d) $f'$ is monotonically increasing.
Put
\begin{align*}
g(x)=\frac{f(x)}{x}\quad(x>0)
\end{align*}
and prove that $g$ is monotonically increasing.

\smallskip
\noindent\textbf{Solution.}\quad

Since $g'(x)=\frac{xf'(x)-f(x)}{x^2}$ where $x>0$, by Theorem 5.11(a), it suffices to show that 
\begin{align*}
xf'(x)-f(x)\ge 0
\end{align*}
for all $x>0$. Since conditions (a), (b), and (c) hold, by the mean value theorem,
\begin{align*}
f(x)=f(x)-f(0)=f'(c)(x-0)=xf'(c)
\end{align*}
for some $c\in(0,x)$. Since $c<x$ here, by condition (d), $xf'(c)\le xf'(x)$. So we conclude that $f(x)\le xf'(x)$.



\bigskip\noindent\hrule\bigskip

Suppose $f'(x)$, $g'(x)$ exist, $g'(x)\ne 0$, and $f(x)=g(x)=0$. Prove that
\begin{align*}
\lim_{t\to x}\frac{f(t)}{g(t)}=\frac{f'(x)}{g'(x)}
\end{align*}
(This holds also for complex functions.)

\smallskip
\noindent\textbf{Solution.}\quad

We have
\begin{align*}
\lim_{t\to x}\frac{f(t)}{g(t)}
=\lim_{t\to x}\frac{\,\dfrac{f(t)-f(x)}{t-x}\,}{\,\dfrac{g(t)-g(x)}{t-x}\,} 
=\frac{\,\displaystyle\lim_{t\to x}\dfrac{f(t)-f(x)}{t-x}\,}{\,\displaystyle\lim_{t\to x}\dfrac{g(t)-g(x)}{t-x}\,} 
=\frac{f'(x)}{g'(x)}
\end{align*}
where the first equality shown above follows by $f(x)=g(x)=0$.



\bigskip\noindent\hrule\bigskip

Suppose $f'$ is continuous on $[a,b]$ and $\varepsilon>0$. Prove that there exists $\delta>0$ such that
\begin{align*}
\left\vert\frac{f(t)-f(x)}{t-x}-f'(x)\right\vert<\varepsilon
\end{align*}
whenever $0<\left\vert t-x\right\vert<\delta$, $a\le x\le b$, $a\le t\le b$. (This could be expressed by saying that $f$ is \textit{uniformly differentiable} on $[a,b]$ if $f'$ is continuous on $[a,b]$.) Does this hold for vector-valued functions too?

\smallskip
\noindent\textbf{Solution.}\quad

(i) By the mean-valued theorem, there exists $c(t)$ between $x$ and $t$ such that $\frac{f(t)-f(x)}{t-x}=f'(c(t))$. Since $f'$ is continuous on $[a,b]$, and since $c(t)\to x$ as $t\to x$ (by the \textit{squeeze theorem}), we have
\begin{align*}
\frac{f(t)-f(x)}{t-x}=f'(c(t))\to f'(x)
\end{align*}
as $t\to x$. This completes the proof.
(ii) Yes. Provided that each component of a given vector-valued function has continuous derivative on an interval.



\bigskip\noindent\hrule\bigskip

Let $f$ be a continuous real function on $\mathbb{R}$, of which it is known that $f'(x)$ exists for all $x\ne 0$ and that $f'(x)\to 3$ as $x\to 0$. Does it follow that $f'(0)$ exists?

\smallskip
\noindent\textbf{Solution.}\quad

Yes. Note that $f'(x)\to 3$ as $x\to 0$, so by L'Hospital's rule,
\begin{align*}
f'(0)
=\lim_{x\to 0}\frac{f(x)-f(0)}{x}
=\lim_{x\to 0}f'(x)
=3
\end{align*}



\bigskip\noindent\hrule\bigskip

Suppose $f$ and $g$ are complex differentiable functions on $(0,1)$, $f(x)\to 0$, $g(x)\to 0$, $f'(x)\to A$, $g'(x)\to B$ as $x\to 0$, where $A$ and $B$ are complex numbers, $B\ne 0$. Prove that
\begin{align*}
\lim_{x\to 0}\frac{f(x)}{g(x)}=\frac{A}{B}
\end{align*}
Compare with Example 5.18. \textit{Hint:}
\begin{align*}
\frac{f(x)}{g(x)}=\left\{\frac{f(x)}{x}-A\right\}\cdot\frac{x}{g(x)}+A\cdot\frac{x}{g(x)}
\end{align*}
Apply Theorem 5.13 to the real and imaginary parts of $f(x)/x$ and $g(x)/x$.

\smallskip
\noindent\textbf{Solution.}\quad

Following the hint, write $f=f_1+if_2$, where $f_1$ and $f_2$ are real functions on $(0,1)$. Then $f_1(x)\to 0$ and $f_1'(x)\to\operatorname{Re}(A)$ as $x\to 0$, and by L'Hospital's rule,
\begin{align*}
\lim_{x\to 0}\operatorname{Re}\left[\frac{f(x)}{x}\right]
=\lim_{x\to 0}\frac{f_1(x)}{x}
=\lim_{x\to 0}f_1'(x)
=\operatorname{Re}(A)
\end{align*}
Similarly we have $\operatorname{Im}\left[\frac{f(x)}{x}\right]\to\operatorname{Im}(A)$, $\operatorname{Re}\left[\frac{g(x)}{x}\right]\to\operatorname{Re}(B)$, and $\operatorname{Im}\left[\frac{g(x)}{x}\right]\to\operatorname{Im}(B)$ as $x\to 0$. So
\begin{align*}
\lim_{x\to 0}\frac{f(x)}{x}
&=\lim_{x\to 0}\left(\operatorname{Re}\left[\frac{f(x)}{x}\right]+i\operatorname{Im}\left[\frac{f(x)}{x}\right]\right) \\
&=\lim_{x\to 0}\operatorname{Re}\left[\frac{f(x)}{x}\right]+i\lim_{x\to 0}\operatorname{Im}\left[\frac{f(x)}{x}\right] \\
&=\operatorname{Re}(A)+i\operatorname{Im}(A) \\
&=A
\end{align*}
and similarly $\lim_{x\to 0}\frac{g(x)}{x}=B$ or the limit of its reciprocal
\begin{align*}
\lim_{x\to 0}\frac{x}{g(x)}
=\lim_{x\to 0}\frac{1}{\,\dfrac{g(x)}{x}\,}
=\frac{1}{\displaystyle\lim_{x\to 0}\dfrac{g(x)}{x}}
=\frac{1}{B}
\end{align*}
exists, since $B\ne 0$. The result will therefore follow.



\bigskip\noindent\hrule\bigskip

Suppose $f$ is defined in a neighborhood of $x$, and suppose $f''(x)$ exists. Show that
\begin{align*}
\lim_{h\to 0}\frac{f(x+h)+f(x-h)-2f(x)}{h^2}=f''(x)
\end{align*}
Show by an example that the limit may exist even $f''(x)$ does not.
\textit{Hint:} Use Theorem 5.13.

\smallskip
\noindent\textbf{Solution.}\quad

(i) By L'Hospital's rule,
\begin{align*}
&\lim_{h\to 0}\frac{f(x+h)+f(x-h)-2f(x)}{h^2} \\
&\qquad=\lim_{h\to 0}\frac{f'(x+h)-f'(x-h)}{2h}\\
&\qquad=\frac{1}{2}\lim_{h\to 0}\left[\frac{f'(x+h)-f'(x)}{h}+\frac{f'(x)-f'(x-h)}{h}\right] \\
&\qquad=\frac{1}{2}\left[\lim_{h\to 0}\frac{f'(x+h)-f'(x)}{h}+\lim_{h\to 0}\frac{f'(x)-f'(x-h)}{h}\right] \\
&\qquad=\frac{1}{2}\left[f''(x)+f''(x)\right] \\
&\qquad=f''(x)
\end{align*}
(ii) For example, let
\begin{align*}
f(x)=\left\{\begin{array}{ll}
x+1&x<0 \vspace{0.08in}\\
0&x=0 \vspace{0.08in}\\
x-1&x>0
\end{array}\right.
\end{align*}



\bigskip\noindent\hrule\bigskip

If $f(x)=\left\vert x\right\vert^3$, compute $f'(x)$, $f''(x)$ for all real $x$, and show that $f^{(3)}(0)$ does not exist.

\smallskip
\noindent\textbf{Solution.}\quad

(i) For $x>0$, $f'(x)=\left(x^3\right)'=3x^2$. For $x<0$, $f'(x)=\left(-x^3\right)'=-3x^2$. Finally,
\begin{align*}
f'(0)
=\lim_{h\to 0}\frac{\left\vert h\right\vert^3}{h}
=\lim_{h\to 0}\operatorname{sgn}(h)h^2
=0
\end{align*}
where
\begin{align*}
\operatorname{sgn}(h)
=\left\{\begin{array}{ll}
1&h>0 \vspace{0.08in}\\
-1&h<0
\end{array}\right.
\end{align*}
(ii) For $x>0$, $f''(x)=\left(3x^2\right)'=6x$. For $x<0$, $f''(x)=\left(-3x^2\right)'=-6x$. Finally, use part (i),
\begin{align*}
f''(0)
=\lim_{h\to 0}\frac{f'(h)}{h}
=\lim_{h\to 0}\frac{\operatorname{sgn}(h)3h^2}{h}
=\operatorname{sgn}(h)3\lim_{h\to 0}h
=0
\end{align*}
(iii) To show that $f^{(3)}(0)$ does not exist, just write
\begin{align*}
\frac{f''(h)-f''(0)}{h}
=\frac{f''(h)}{h}
=\frac{\operatorname{sgn}(h)6h}{h}
=6\operatorname{sgn}(h)
\end{align*}
The result then follows.



\bigskip\noindent\hrule\bigskip

Suppose $a$ and $x$ are real numbers, $c>0$, and $f$ is defined on $[-1,1]$ by
\begin{align*}
f(x)=\left\{\begin{array}{ll}
x^a\sin\left(\left\vert x\right\vert^{-c}\right)&\mbox{(if $x\ne 0$)}\vspace{0.08in}\\
0&\mbox{(if $x=0$)}
\end{array}\right.
\end{align*}
Prove the following statements:
(a) $f$ is continuous if and only if $a>0$.
(b) $f'(0)$ exists if and only if $a>1$.
(c) $f'$ is bounded if and only if $a\ge 1+c$.
(d) $f'$ is continuous if and only if $a>1+c$.
(e) $f''(0)$ exists if and only if $a>2+c$.
(f) $f''$ is bounded if and only if $a\ge 2+2c$.
(g) $f''$ is continuous if and only if $a>2+2c$.

\smallskip
\noindent\textbf{Solution.}\quad

(a) It is clear that $f$ is continuous at any nonzero point $x$, so we only consider the point $x=0$. If $a>0$, then
\begin{align*}
\left\vert f(x)\right\vert
=\left\vert x^a\sin\left(\left\vert x\right\vert^{-c}\right)\right\vert
\le \left\vert x^a\right\vert
\to 0\quad\mbox{as }x\to 0
\end{align*}
By squeezing, $f(x)\to 0$ as $x\to 0$, and hence $f$ is continuous at $0$. Conversely, if $a\le 0$, let $x_n=\left(2n\pi+\frac{\pi}{2}\right)^{-\frac{1}{c}}$ for $n=1,2,\ldots$. Then $x_n\to 0$ as $n\to\infty$, but $-\frac{a}{c}\ge 0$ implies
\begin{align*}
f(x_n)=\left(2n\pi+\frac{\pi}{2}\right)^{-\frac{a}{c}}\ge 1\quad\mbox{for all }n
\end{align*}
This shows that $\lim_{n\to\infty}f(x_n)\ne 0=f(0)$, and hence $f$ is not continuous at $0$.

\smallskip
\noindent\textbf{Solution.}\quad
(b) If $a>1$ or $a-1>0$, then
\begin{align*}
\left\vert\frac{f(x)-f(0)}{x}\right\vert
=\left\vert\frac{x^a\sin\left(\left\vert x\right\vert^{-c}\right)}{x}\right\vert
\le\left\vert x^{a-1}\right\vert
\to 0\quad\mbox{as }x\to 0
\end{align*}
By squeezing, $f'(0)=0$. Conversely, if $a\le 1$ or $a-1\le 0$, let $x_n=\left(2n\pi+\frac{\pi}{2}\right)^{-\frac{1}{c}}$ and $y_n=(2n\pi)^{-\frac{1}{c}}$ for $n=1,2,\ldots$. Then $x_n\to 0$ and $y_n\to 0$ as $n\to\infty$, but $-\frac{a-1}{c}\ge 0$ implies
\begin{align*}
\frac{f(x_n)-f(0)}{x_n}
&=\left(2n\pi+\frac{\pi}{2}\right)^{-\frac{a-1}{c}}
\ge 1\quad\mbox{for all }n
\end{align*}
Also, we have $\frac{f(y_n)-f(0)}{y_n}=0$ for all $n$. This gives
\begin{align*}
\lim_{n\to\infty}\frac{f(x_n)-f(0)}{x_n}\ne 0=\lim_{n\to\infty}\frac{f(y_n)-f(0)}{y_n}
\end{align*}
and hence $f'(0)$ does not exist.

\smallskip
\noindent\textbf{Solution.}\quad
(c)

\smallskip
\noindent\textbf{Solution.}\quad
(d)

\smallskip
\noindent\textbf{Solution.}\quad
(e)

\smallskip
\noindent\textbf{Solution.}\quad
(f)

\smallskip
\noindent\textbf{Solution.}\quad
(g)



\bigskip\noindent\hrule\bigskip

Let $f$ be a differentiable real function defined on $(a,b)$. Prove that $f$ is convex if and only if $f'$ is monotonically increasing. Assume next that $f''(x)$ exists for every $x\in(a,b)$, and prove that $f$ is convex if and only if $f''(x)\ge 0$ for all $x\in(a,b)$.

\smallskip
\noindent\textbf{Solution.}\quad

(i) If $f$ is convex. For $x,y,t\in(a,b)$ with $x<t<y$, and let $c=\frac{f(y)-f(x)}{y-x}$. By Exercise 4.23,
\begin{align*}
c\ge\frac{f(t)-f(x)}{t-x}\quad\mbox{and}\quad
c\le\frac{f(y)-f(t)}{y-t}
\end{align*}
Since $f$ is differentiable, we have $c\ge f'(x)$ and $c\le f'(y)$ as $t\to x$ and $t\to y$, respectively. Hence $f'(x)\le f'(y)$. Conversely, if $f'$ is monotonically increasing, given $x,y\in(a,b)$ with $x<y$, and $\lambda\in(0,1)$. Put $t=\lambda x+(1-\lambda)y$, then
\begin{align*}
t-x=(1-\lambda)(y-x)\quad\mbox{and}\quad
y-t=\lambda(y-x)
\end{align*}
Since $f'$ is monotonically increasing, applying the mean value theorem,
\begin{align*}
\frac{f(t)-f(x)}{t-x}=f'(c_1)\le f'(c_2)=\frac{f(y)-f(t)}{y-t}
\end{align*}
for some $c_1\in(x,t)$ and $c_2\in(t,y)$. It follows that
\begin{align*}
\lambda\left[f(t)-f(x)\right]
\le(1-\lambda)\left[f(y)-f(t)\right]
\end{align*}
or
\begin{align*}
f(t)
\le\lambda f(x)+(1-\lambda)f(y)
\end{align*}
(ii) We first show that $f'$ is monotonically increasing if and only if $f''(x)\ge 0$ for all $x\in(a,b)$. The sufficiency part is the result of Theorem 5.11(a). To show the necessity part, given $x\in(a,b)$, observe that the fraction $\frac{f'(t)-f'(x)}{t-x}$ is always nonnegative, for all $t\in(a,b)$ with $t\ne x$, since $f'$ is monotonically increasing. Thus by definition,
\begin{align*}
f''(x)=\lim_{t\to x}\frac{f'(t)-f'(x)}{t-x}\ge 0
\end{align*}
Hence, $f$ is convex if and only if $f'$ is monotonically increasing [by part (i)], if and only if $f''(x)\ge 0$ for all $x\in(a,b)$.



\bigskip\noindent\hrule\bigskip

Suppose $a\in \mathbb{R}$, $f$ is a twice-differentiable real function on $(a,\infty)$, and $M_0$, $M_1$, $M_2$ are the least upper bounds of $\left\vert f(x)\right\vert$, $\left\vert f'(x)\right\vert$, $\left\vert f''(x)\right\vert$, respectively, on $(a,\infty)$. Prove that
\begin{align*}
M_1^2\le 4M_0M_2
\end{align*}
\textit{Hint:} If $h>0$, Taylor's theorem shows that
\begin{align*}
f'(x)=\frac{1}{2h}\left[f(x+2h)-f(x)\right]-hf''(\xi)
\end{align*}
for some $\xi\in(x,x+2h)$. Hence
\begin{align*}
\left\vert f'(x)\right\vert\le hM_2+\frac{M_0}{h}
\end{align*}
To show that $M_1^2=4M_0M_2$ can actually happen, take $a=-1$, define
\begin{align*}
f(x)=\left\{\begin{array}{ll}
2x^2-1&\mbox{($-1<x<0$)}\vspace{0.1in}\\
\dfrac{x^2-1}{x^2+1}&\mbox{($0\le x<\infty$)}
\end{array}\right.
\end{align*}
and show that $M_0=1$, $M_1=4$, $M_2=4$.
Does $M_1^2\le 4M_0M_2$ hold for vector-valued functions too?

\smallskip
\noindent\textbf{Solution.}\quad

(i) First, note that $M_0M_2$ be of the form $0\cdot\infty$ is exceptional. Now we consider the following four cases.
Case 1: $M_0=\infty$ or $M_2=\infty$. The result is trivial.
Case 2: $M_0=0$. Then $f(x)=0$, and then $f'(x)=f''(x)=0$ for all $x\in(a,\infty)$. It follows that $M_1=M_2=0$, and the result is also trivial.
Case 3: $0<M_0<\infty$ and $M_2=0$. Then $f''(x)=0$, and then $f'(x)=c$ and $f(x)=cx+d$ for some constants $c,d\in R$, for all $x\in(a,\infty)$. Since $M_0$ is finite, we need $c=0$ and then $M_1=0$. The result therefore follows.
Case 4: $0<M_0<\infty$ and $0<M_2<\infty$. Following the hint, use Taylor's theorem (Theorem 5.15), for $x\in(a,\infty)$ and $h>0$ there exists $\xi\in(x,x+2h)$ such that
\begin{align*}
f(x+2h)
&=f(x)+f'(x)\cdot 2h+\frac{f''(\xi)}{2}\cdot (2h)^2 \\
&=f(x)+f'(x)\cdot 2h+f''(\xi)\cdot 2h^2
\end{align*}
Thus
\begin{align*}
f'(x)=\frac{1}{2h}\left[f(x+2h)-f(x)\right]-hf''(\xi)
\end{align*}
and hence
\begin{align*}
\left\vert f'(x)\right\vert\le hM_2+\frac{M_0}{h}
\end{align*}
By assumption, we may put $h=\sqrt{\frac{M_0}{M_2}}$ and the inequality becomes
\begin{align*}
\left\vert f'(x)\right\vert\le 2\sqrt{M_0M_2}
\end{align*}
implying $M_1\le 2\sqrt{M_0M_2}$ or $M_1^2\le 4M_0M_2$.
(ii) By some calculation,
\begin{align*}
(-1<x<0)&\left\{\begin{array}{l}
f'(x)=4x \vspace{0.08in}\\
f''(x)=4
\end{array}\right. \\
(0<x<\infty)&\left\{\begin{array}{l}
f'(x)=\dfrac{4x}{(x^2+1)^2} \vspace{0.08in}\\
f''(x)=\dfrac{-4(3x^2-1)}{(x^2+1)^3}
\end{array}\right.
\end{align*}
To find $f'(0)$ and $f''(0)$, calculate
\begin{align*}
\lim_{t\to 0+}\frac{f(t)-f(0)}{t}
&=\lim_{t\to 0+}\frac{\dfrac{t^2-1}{t^2+1}+1}{t}
=\lim_{t\to 0+}\frac{2t}{t^2+1}
=0 \\
\lim_{t\to 0-}\frac{f(t)-f(0)}{t}
&=\lim_{t\to 0-}\frac{(2t^2-1)+1}{t}
=\lim_{t\to 0-}2t
=0
\end{align*}
Then $f'(0)=0$. Next,
\begin{align*}
\lim_{t\to 0+}\frac{f'(t)-f'(0)}{t}
&=\lim_{t\to 0+}\frac{\,\dfrac{4t}{(t^2+1)^2}\,}{t}
=\lim_{t\to 0+}\frac{4}{(t^2+1)^2}
=4 \\
\lim_{t\to 0-}\frac{f'(t)-f'(0)}{t}
&=\lim_{t\to 0-}\frac{4t}{t}
=4
\end{align*}
Then $f''(0)=4$. Now, $-1\le f(x)<1$ for $x\in(-1,\infty)$, and $\lim_{x\to\infty}f(x)=1$, are both trivial. Hence $M_0=1$. Also, since $f''(x)=4$ for $x\in(-1,0]$, and $\lim_{x\to 0+}f''(x)=4$, we see that $f$ is twice-differentiable on $(-1,\infty)$. To show that $\left\vert f''(x)\right\vert<4$ for all $x>0$, observe that
\begin{align*}
3x^2-1<3x^2+1<x^6+3x^4+3x^2+1=(x^2+1)^3
\end{align*}
and that $\left\vert\frac{3x^2-1}{(x^2+1)^3}\right\vert=\frac{3x^2-1}{(x^2+1)^3}<1$ for $x>0$. So
\begin{align*}
\left\vert f''(x)\right\vert=4\left\vert\frac{3x^2-1}{(x^2+1)^3}\right\vert<4
\end{align*}
for $x>0$, and hence $M_2=4$. Finally, since $0\le f'(x)<4$ for $x\in(-1,0]$, and
\begin{align*}
2x\le 1+x^2<1+2x^2+x^4=(x^2+1)^2
\end{align*}
for $x>0$. We see that 
\begin{align*}
\left\vert f'(x)\right\vert=\left\vert\frac{4x}{(x^2+1)^2}\right\vert=2\left[\frac{2x}{(x^2+1)^2}\right]<2<4
\end{align*}
for $x>0$, and hence $M_1=4$ (since $\lim_{t\to(-1)+}f'(x)=-4$).
(iii) Yes. In $\mathbb{R}^k$, let $\mathbf{f}(x)=(f_1(x),\ldots,f_k(x))$ where each $f_j$ ($1\le j\le k$) is a twice-differentiable real function. Note that $M_0M_2$ be of the form $0\cdot\infty$ is exceptional, and we consider the following four cases. (The first three cases are simple extensions of $k=1$.)
Case 1: $M_0=\infty$ or $M_2=\infty$. The result is trivial.
Case 2: $M_0=0$. Then $f_j(x)=0$, and then $f_j'(x)=f_j''(x)=0$ for all $x\in(a,\infty)$ and for all $j$. It follows that $M_1=M_2=0$, and the result is also trivial.
Case 3: $0<M_0<\infty$ and $M_2=0$. Then $f_j''(x)=0$, and then $f_j'(x)=c_j$ and $f_j(x)=c_jx+d_j$ for some constants $c_j,d_j\in R$, for all $x\in(a,\infty)$ and for all $j$. Since $M_0$ is finite, we need $c_j=0$ for each $j$, and then $M_1=0$. The result therefore follows.
Case 4: $0<M_0<\infty$ and $0<M_2<\infty$. If $M_1=0$ then we are done. If $M_1>0$, let $p\in \mathbb{R}$ be such that $0<p<M_1$, and let $x_0\in(a,\infty)$ be such that $\left\vert\mathbf{f}'(x_0)\right\vert>p$. Put $\mathbf{u}=\frac{\mathbf{f}'(x_0)}{\left\vert\mathbf{f}'(x_0)\right\vert}$. Consider the real-valued function $g(x)=\mathbf{u}\cdot\mathbf{f}(x)$ for $x\in(a,\infty)$, and note that $g$ is twice-differentiable. Let $N_0$, $N_1$, $N_2$ be the least upper bounds of $\left\vert g(x)\right\vert$, $\left\vert g'(x)\right\vert$, $\left\vert g''(x)\right\vert$, respectively. Since $\left\vert\mathbf{u}\right\vert=1$, by Schwarz inequality [Theorem 1.37(d)],
\begin{align*}
\left\vert g(x)\right\vert
\le\left\vert\mathbf{u}\right\vert\left\vert\mathbf{f}(x)\right\vert
=\left\vert\mathbf{f}(x)\right\vert,\quad
\left\vert g''(x)\right\vert
\le\left\vert\mathbf{u}\right\vert\left\vert\mathbf{f}''(x)\right\vert
=\left\vert\mathbf{f}''(x)\right\vert
\end{align*}
for all $x\in(a,\infty)$. So that $N_0\le M_0$ and $N_2\le M_2$. Also, since
\begin{align*}
N_1\ge g'(x_0)=\mathbf{u}\cdot\mathbf{f}'(x_0)=\left\vert\mathbf{f}'(x_0)\right\vert>p
\end{align*}
and since $N_1\le 4N_0N_2$ [by part (i)], we have $p<4M_0M_2$. Since $p$ is arbitrarily chosen such that $0<p<M_1$, we conclude that $M_1\le 4M_0M_2$.



\bigskip\noindent\hrule\bigskip

Suppose $f$ is twice-differentiable on $(0,\infty)$, $f''$ is bounded on $(0,\infty)$, and $f(x)\to 0$ as $x\to\infty$. Prove that $f'(x)\to 0$ as $x\to\infty$.
\textit{Hint:} Let $a\to\infty$ in Exercise 15.

\smallskip
\noindent\textbf{Solution.}\quad

Let $M>0$ be such that $\left\vert f''(x)\right\vert\le M$ for $x\in(0,\infty)$. Given $a\in(0,\infty)$, let $M_0(a)$, $M_1(a)$, $M_2(a)$ be the least upper bounds of $\left\vert f(x)\right\vert$, $\left\vert f'(x)\right\vert$, $\left\vert f''(x)\right\vert$, respectively, for $x\in(a,\infty)$. Then clearly $M_2(a)\le M$ for all $a$. Since $f(x)\to 0$ as $x\to\infty$, for every $\varepsilon>0$ there exists $a_0\in(0,\infty)$ such that $\left\vert f(x)\right\vert<\varepsilon$ for all $x\in(a_0,\infty)$. It follows that $M_0(a_0)\le\varepsilon$ and that (by Exercise 15)
\begin{align*}
M_1^2(a_0)\le 4M_0(a_0)M_2(a_0)\le 4M\varepsilon
\end{align*}
i.e., $\left\vert f'(x)\right\vert\le 2\sqrt{M\varepsilon}$ for all $x\in(a_0,\infty)$. Hence $f'(x)\to 0$ as $x\to\infty$.



\bigskip\noindent\hrule\bigskip

Suppose $f$ is a real, three times differentiable function on $[-1,1]$, such that
\begin{align*}
f(-1)=0,\quad f(0)=0,\quad f(1)=1,\quad f'(0)=0
\end{align*}
Prove that $f^{(3)}(x)\ge 3$ for some $x\in(-1,1)$.
Note that equality holds for $\tfrac{1}{2}(x^3+x^2)$.
\textit{Hint:} Use Theorem 5.15, with $\alpha=0$ and $\beta=\pm 1$, to show that there exist $s\in(0,1)$ and $t\in(-1,0)$ such that
\begin{align*}
f^{(3)}(s)+f^{(3)}(t)=6
\end{align*}

\smallskip
\noindent\textbf{Solution.}\quad

Following the hint, use Taylor's theorem (Theorem 5.15), there exists $s\in(0,1)$ and $t\in(-1,0)$ such that
\begin{align*}
1
&=f(1)
=f(0)+f'(0)+\frac{f''(0)}{2}+\frac{f^{(3)}(s)}{6}
=\frac{f''(0)}{2}+\frac{f^{(3)}(s)}{6} \\
0
&=f(-1)
=f(0)-f'(0)+\frac{f''(0)}{2}-\frac{f^{(3)}(s)}{6}
=\frac{f''(0)}{2}-\frac{f^{(3)}(t)}{6}
\end{align*}
Subtract them we get
\begin{align*}
f^{(3)}(s)+f^{(3)}(t)=6
\end{align*}
It follows that either $f^{(3)}(s)\ge 3$ or $f^{(3)}(t)\ge 3$, which completes the proof.



\bigskip\noindent\hrule\bigskip

Suppose $f$ is a real function on $[a,b]$, $n$ is a positive integer, and $f^{(n-1)}$ exists for every $t\in[a,b]$. Let $\alpha$, $\beta$, and $P$ be as in Taylor's theorem (5.15). Define
\begin{align*}
Q(t)=\frac{f(t)-f(\beta)}{t-\beta}
\end{align*}
for $t\in[a,b]$, $t\ne\beta$, differentiate
\begin{align*}
f(t)-f(\beta)=(t-\beta)Q(t)
\end{align*}
$n-1$ times at $t=\alpha$, and derive the following version of Taylor's theorem:
\begin{align*}
f(\beta)=P(\beta)+\frac{Q^{(n-1)}(\alpha)}{(n-1)!}(\beta-\alpha)^n
\end{align*}

\smallskip
\noindent\textbf{Solution.}\quad

By induction, we easily conclude that
\begin{align*}
f^{(k)}(\alpha)=kQ^{(k-1)}(\alpha)-(\beta-\alpha)Q^{(k)}(\alpha)
\end{align*}
and then
\begin{align*}
\frac{f^{(k)}(\alpha)}{k!}\left(\beta-\alpha\right)^k
=\frac{Q^{(k-1)}(\alpha)}{(k-1)!}\left(\beta-\alpha\right)^k-\frac{Q^{(k)}(\alpha)}{k!}\left(\beta-\alpha\right)^{k+1}
\end{align*}
for $k=1,2,\ldots,n-1$. It follows that
\begin{align*}
P(\beta)
&=f(\alpha)+\sum_{k=1}^{n-1}\frac{f^{(k)}(\alpha)}{k!}\left(\beta-\alpha\right)^k \\
&=f(\alpha)+Q(\alpha)(\beta-\alpha)-\frac{Q^{(n-1)}(\alpha)}{(n-1)!}\left(\beta-\alpha\right)^n \\
&=f(\alpha)+\left[f(\beta)-f(\alpha)\right]-\frac{Q^{(n-1)}(\alpha)}{(n-1)!}\left(\beta-\alpha\right)^n \\
&=f(\beta)-\frac{Q^{(n-1)}(\alpha)}{(n-1)!}\left(\beta-\alpha\right)^n \\
\end{align*}
or $f(\beta)=P(\beta)+\frac{Q^{(n-1)}(\alpha)}{(n-1)!}\left(\beta-\alpha\right)^n$.



\bigskip\noindent\hrule\bigskip

Suppose $f$ is defined in $(-1,1)$ and $f'(0)$ exists. Suppose $-1<\alpha_n<\beta_n<1$, $\alpha_n\to 0$, and $\beta_n\to 0$ as $n\to\infty$. Define the difference quotients
\begin{align*}
D_n=\frac{f(\beta_n)-f(\alpha_n)}{\beta_n-\alpha_n}
\end{align*}
Prove the following statements:
(a) If $\alpha_n<0<\beta_n$, then $\lim D_n=f'(0)$.
(b) If $0<\alpha_n<\beta_n$ and $\{\beta_n/(\beta_n-\alpha_n)\}$ is bounded, then $\lim D_n=f'(0)$.
(c) If $f'$ is continuous in $(-1,1)$, then $\lim D_n=f'(0)$.
Give an example in which $f$ is differentiable in $(-1,1)$ (but $f'$ is not continuous at $0$) and in which $\alpha_n$, $\beta_n$ tend to $0$ in such a way that $\lim D_n$ exists but is different from $f'(0)$.

\smallskip
\noindent\textbf{Solution.}\quad

(a) Denote $\lambda_n=\frac{\beta_n}{\beta_n-\alpha_n}$ for $n=1,2,\ldots$, then clearly $0<\lambda_n<1$ (since $\alpha_n<0<\beta_n$). Now, we can write
\begin{align*}
D_n
&=\frac{f(\beta_n)-f(0)+f(0)-f(\alpha_n)}{\beta_n-\alpha_n} \\
&=\frac{\beta_n}{\beta_n-\alpha_n}\cdot\frac{f(\beta_n)-f(0)}{\beta_n}
+\frac{-\alpha_n}{\beta_n-\alpha_n}\cdot\frac{f(0)-f(\alpha_n)}{-\alpha_n} \\
&=\lambda_n\cdot\frac{f(\beta_n)-f(0)}{\beta_n}
+\left(1-\lambda_n\right)\cdot\frac{f(\alpha_n)-f(0)}{\alpha_n}
\end{align*}
Since $f'(0)$ exists, by definition, for $\varepsilon>0$ there exists a positive integer $N$ such that $n\ge N$ implies
\begin{align*}
\left\vert\frac{f(\beta_n)-f(0)}{\beta_n}-f'(0)\right\vert<\varepsilon\quad\mbox{and}\quad
\left\vert\frac{f(\alpha_n)-f(0)}{\alpha_n}-f'(0)\right\vert<\varepsilon
\end{align*}
So $n\ge N$ implies
\begin{align*}
\left\vert D_n-f'(0)\right\vert
&=\left\vert \lambda_n\left[\frac{f(\beta_n)-f(0)}{\beta_n}-f'(0)\right]
+(1-\lambda_n)\left[\frac{f(\alpha_n)-f(0)}{\alpha_n}-f'(0)\right]\right\vert \\
&\le\lambda_n\left\vert\frac{f(\beta_n)-f(0)}{\beta_n}-f'(0)\right\vert
+(1-\lambda_n)\left\vert\frac{f(\alpha_n)-f(0)}{\alpha_n}-f'(0)\right\vert \\
&<\lambda_n\varepsilon+\left(1-\lambda_n\right)\varepsilon \\
&=\varepsilon
\end{align*}
Hence $\lim_{n\to\infty}D_n=f'(0)$.

\smallskip
\noindent\textbf{Solution.}\quad
(b) Let $\lambda_n$ be the same notation as in part (a), then we know that $\lambda_n>0$ for each $n$. Since $\{\beta_n/(\beta_n-\alpha_n)\}$ is bounded, there exists $M>0$ such that $\lambda_n\le M$ for each $n$. Now, we can write
\begin{align*}
D_n
&=\frac{\left[f(\beta_n)-f(0)\right]-\left[f(\alpha_n)-f(0)\right]}{\beta_n-\alpha_n} \\
&=\frac{\beta_n}{\beta_n-\alpha_n}\cdot\frac{f(\beta_n)-f(0)}{\beta_n}
+\frac{-\alpha_n}{\beta_n-\alpha_n}\cdot\frac{f(\alpha_n)-f(0)}{\alpha_n} \\
&=\lambda_n\cdot\frac{f(\beta_n)-f(0)}{\beta_n}
+(1-\lambda_n)\cdot\frac{f(\alpha_n)-f(0)}{\alpha_n}
\end{align*}
Then if we let $\varepsilon$ and $N$ be the same meaning as in part (a), it follows that $n\ge N$ implies
\begin{align*}
\left\vert D_n-f'(0)\right\vert
&=\left\vert \lambda_n\left[\frac{f(\beta_n)-f(0)}{\beta_n}-f'(0)\right]
+(1-\lambda_n)\left[\frac{f(\alpha_n)-f(0)}{\alpha_n}-f'(0)\right]\right\vert \\
&\le\lambda_n\left\vert\frac{f(\beta_n)-f(0)}{\beta_n}-f'(0)\right\vert
+(1+\lambda_n)\left\vert\frac{f(\alpha_n)-f(0)}{\alpha_n}-f'(0)\right\vert \\
&<\lambda_n\varepsilon+\left(1+\lambda_n\right)\varepsilon \\
&\le\left(1+2M\right)\varepsilon
\end{align*}
Hence $\lim_{n\to\infty}D_n=f'(0)$.

\smallskip
\noindent\textbf{Solution.}\quad
(c) By the mean value theorem, for each $n$ there exists $\gamma_n\in(\alpha_n,\beta_n)$ such that
\begin{align*}
D_n=f'(\gamma_n)
\end{align*}
By squeezing, we see that $\gamma_n\to 0$ as $n\to\infty$. Since $f'$ is continuous, $D_n=f'(\gamma_n)\to f'(0)$ as $n\to\infty$.

\smallskip
\noindent\textbf{Solution.}\quad
Inspired from Exercise 13, we give an example as below. Let $f:(-1,1)\to \mathbb{R}$ be defined by
\begin{align*}
f(x)=\left\{\begin{array}{ll}
0&x=0 \vspace{0.08in}\\
x^2\sin\left(\dfrac{1}{x}\right)&\mbox{otherwise}
\end{array}\right.
\end{align*}
Then $f'(x)=2x\sin\left(\frac{1}{x}\right)-\cos\left(\frac{1}{x}\right)$ for $x\ne0$, and
\begin{align*}
f'(0)
=\lim_{x\to 0}\frac{f(x)-f(0)}{x}
=\lim_{x\to 0}x\sin\left(\frac{1}{x}\right)
=0
\end{align*}
i.e., such $f$ is differentiable on $(-1,1)$ but $f'$ is not continuous at $0$. Moreover, let $\alpha_n=\left(2n\pi+\frac{\pi}{2}\right)^{-1}$ and $\beta_n=\left(2n\pi\right)^{-1}$ for $n=1,2,\ldots$. Then $0<\alpha_n<\beta_n<1$, $\alpha_n\to 0$, and $\beta_n\to 0$ as $n\to\infty$. But
\begin{align*}
D_n
&=\frac{-\left(2n\pi+\dfrac{\pi}{2}\right)^{-2}}{\left(2n\pi\right)^{-1}-\left(2n\pi+\dfrac{\pi}{2}\right)^{-1}} \\
&=\frac{-1}{\left(2n\pi+\dfrac{\pi}{2}\right)^{2}\left(2n\pi\right)^{-1}-\left(2n\pi+\dfrac{\pi}{2}\right)} \\
&=\frac{-1}{\left(2n\pi+\dfrac{\pi}{2}\right)\left(1+\dfrac{1}{4n}\right)-\left(2n\pi+\dfrac{\pi}{2}\right)} \\
&=\frac{-4n}{2n\pi+\dfrac{\pi}{2}}
\end{align*}
which follows that $\lim_{n\to\infty}D_n=-\frac{2}{\pi}\ne 0=f'(0)$.



\bigskip\noindent\hrule\bigskip

Formulate and prove an inequality which follows form Taylor's theorem and which remains valid for vector-valued functions.

\smallskip
\noindent\textbf{Solution.}\quad

(i) The statement in the real-valued case: \textit{Suppose $f$ is a real function on $[a,b]$, $n$ is a positive integer, $f^{(n-1)}$ is continuous on $[a,b]$, $f^{(n)}(t)$ exists for every $t\in(a.b)$. Let
\begin{align*}
P(t)=\sum_{k=0}^{n-1}\frac{f^{(k)}(a)}{k!}(t-a)^k
\end{align*}
then there exists $x\in(a,b)$ such that}
\begin{align*}
\left\vert f(b)-P(b)\right\vert\le\left\vert\frac{f^{(n)}(x)}{n!}\right\vert(b-a)^n
\end{align*}
The proof of the assertion directly follows from the Taylor' theorem.
(ii) The statement in the vector-valued case: \textit{Suppose $\mathbf{f}$ is a mapping on $[a,b]$ into $\mathbb{R}^k$, $n$ is a positive integer, $\mathbf{f}^{(n-1)}$ is continuous on $[a,b]$, $\mathbf{f}^{(n)}(t)$ exists for every $t\in(a.b)$. Let
\begin{align*}
\mathbf{P}(t)=\sum_{k=0}^{n-1}\frac{\mathbf{f}^{(k)}(a)}{k!}(t-a)^k
\end{align*}
then there exists $x\in(a,b)$ such that}
\begin{align*}
\left\vert \mathbf{f}(b)-\mathbf{P}(b)\right\vert\le\left\vert\frac{\mathbf{f}^{(n)}(x)}{n!}\right\vert(b-a)^n
\end{align*}
To show it, note that $\left\vert \mathbf{f}(b)-\mathbf{P}(b)\right\vert=0$ is a trivial case. If $\left\vert \mathbf{f}(b)-\mathbf{P}(b)\right\vert\ne 0$, define 
\begin{align*}
\mathbf{u}=\frac{1}{\left\vert \mathbf{f}(b)-\mathbf{P}(b)\right\vert}\left[\mathbf{f}(b)-\mathbf{P}(b)\right]
\end{align*}
Then $\mathbf{u}\cdot\mathbf{f}$ is a real-valued function on $[a,b]$, satisfying all the conditions in the statement of part (i). Since $(\mathbf{u}\cdot\mathbf{f})^{(k)}=\mathbf{u}\cdot\mathbf{f}^{(k)}$ for $k=0,1,2,\ldots,n$, and
\begin{align*}
\sum_{k=0}^{n-1}\frac{(\mathbf{u}\cdot\mathbf{f})^{(k)}(a)}{k!}(t-a)^k
&=\sum_{k=0}^{n-1}\frac{\mathbf{u}\cdot\mathbf{f}^{(k)}(a)}{k!}(t-a)^k \\
&=\mathbf{u}\cdot\left[\sum_{k=0}^{n-1}\frac{\mathbf{f}^{(k)}(a)}{k!}(t-a)^k\right] \\
&=\mathbf{u}\cdot\mathbf{P}(t)
\end{align*}
by part (i) we have
\begin{align*}
\left\vert \mathbf{u}\cdot\mathbf{f}(b)-\mathbf{u}\cdot\mathbf{P}(b)\right\vert
\le\left\vert\frac{\mathbf{u}\cdot\mathbf{f}^{(n)}(x)}{n!}\right\vert(b-a)^n
\le\left\vert\frac{\mathbf{f}^{(n)}(x)}{n!}\right\vert(b-a)^n
\end{align*}
for some $x\in(a,b)$, where the above second inequality follows by the fact $\left\vert\mathbf{u}\right\vert=1$, and by the Schwarz inequality. On the other hand,
\begin{align*}
\mathbf{u}\cdot\mathbf{f}(b)-\mathbf{u}\cdot\mathbf{P}(b)
=\mathbf{u}\cdot\left[\mathbf{f}(b)-\mathbf{P}(b)\right]
=\left\vert\mathbf{f}(b)-\mathbf{P}(b)\right\vert
\end{align*}
we then conclude that
\begin{align*}
\left\vert\mathbf{f}(b)-\mathbf{P}(b)\right\vert
\le\left\vert\frac{\mathbf{f}^{(n)}(x)}{n!}\right\vert(b-a)^n
\end{align*}



\bigskip\noindent\hrule\bigskip

Let $E$ be a closed subset of $\mathbb{R}$. We saw in Exercise 22, Chap. 4, that there is a real continuous function $f$ on $\mathbb{R}$ whose zero set is $E$. It is possible, for each closed set $E$, to find such an $f$ which is differentiable on $\mathbb{R}$, or one which is $n$ times differentiable, or even one which has derivatives of all orders on $\mathbb{R}$?

\smallskip
\noindent\textbf{Solution.}\quad

Yes.  Write the open set $\mathbb{R}\setminus E$ as the disjoint union of its component intervals.  On a bounded component $(a,b)$ put
\[
 h_{a,b}(x)=
 \begin{cases}
 \exp\!\left(-\dfrac{1}{(x-a)^2}-\dfrac{1}{(b-x)^2}\right),&a<x<b,\\
 0,&\text{otherwise}.
 \end{cases}
\]
For a component $(a,\infty)$ use $\exp(-1/(x-a)^2)$ when $x>a$ and zero otherwise; define the analogous function on $(-\infty,b)$.  Every derivative of these functions tends to zero at each finite endpoint.  Define $f$ on each component to be the corresponding positive function and put $f=0$ on $E$.  Since the components are disjoint, this definition is unambiguous.  At an accumulation point of component endpoints, each derivative still tends to zero: on a bounded component of length $\ell$, the displayed exponential and each fixed-order derivative are bounded by a polynomial in $1/\ell$ times $e^{-c/\ell^2}$, which tends to zero faster than every power of $\ell$.  Thus $f\in C^\infty(\mathbb{R})$, it is positive precisely on $\mathbb{R}\setminus E$, and its zero set is exactly $E$.  The cases $E=\mathbb{R}$ and $E=\varnothing$ are given by $f=0$ and $f=1$, respectively.



\bigskip\noindent\hrule\bigskip

Suppose $f$ is a real function on $(-\infty,\infty)$. Call $x$ a \textit{fixed point} of $f$ if $f(x)=x$.
(a) If $f$ is differentiable and $f'(t)\ne 1$ for every real $t$, prove that $f$ has at most one fixed point.
(b) Show that the function $f$ defined by
\begin{align*}
f(t)=t+(1+e^t)^{-1}
\end{align*}
has no fixed point, although $0<f'(t)<1$ for all real $t$.
(c) However, if there is a constant $A<1$ such that $\left\vert f'(t)\right\vert\le A$ for all real $t$, prove that a fixed point $x$ of $f$ exists, and that $x=\lim x_n$, where $x_1$ is an arbitrary real number and
\begin{align*}
x_{n+1}=f(x_n)
\end{align*}
for $x=1,2,3,\ldots$.
(d) Show that the process described in (c) can be visualized by the zig-zag path
\begin{align*}
(x_1,x_2)\to(x_2,x_2)\to(x_2,x_3)\to(x_3,x_3)\to(x_3,x_4)\to\cdots
\end{align*}

\smallskip
\noindent\textbf{Solution.}\quad

(a) Suppose there are two fixed points of $f$, saying $x$ and $y$. W.l.o.g., let $x<y$. Since $f$ is differentiable, by the mean value theorem, there exists a point $t\in(x,y)$ such that $f(y)-f(x)=(y-x)f'(t)$. But since $f(x)=x$ and $f(y)=y$, we have
\begin{align*}
y-x
=(y-x)f'(t)
\end{align*}
implying $f'(t)=1$.

\smallskip
\noindent\textbf{Solution.}\quad
(b) Note that $f(t)=t$ if and only if $(1+e^t)^{-1}=0$, which is impossible. So $f$ has no fixed point. Next, observe that
\begin{align*}
f'(t)=1-\frac{e^t}{(1+e^t)^{2}}\in(0,1)
\end{align*}
since $\frac{e^t}{(1+e^t)^{2}}\in(0,1)$.

\smallskip
\noindent\textbf{Solution.}\quad
(c) The uniqueness of the fixed point $x$ of $f$ follows by part (a). To show the existence, given $x_1\in \mathbb{R}$ and define
\begin{align*}
x_{n+1}=f(x_n)
\end{align*}
for $n=1,2,\ldots$. Then by induction and by the mean value theorem,
\begin{align*}
\left\vert x_{n+1}-x_n\right\vert
&=\left\vert f(x_n)-f(x_{n-1})\right\vert \\
&=\left\vert f'(c_n)\left(x_n-x_{n-1}\right)\right\vert \\
&\le A\left\vert x_n-x_{n-1}\right\vert \\
&\le A^2\left\vert x_{n-1}-x_{n-2}\right\vert \\
&\,\,\vdots \\
&\le A^{n-1}\left\vert x_2-x_1\right\vert
\end{align*}
Since $0<A<1$, for every $\varepsilon>0$ there exists a positive $N$ such that $\frac{A^{N}}{1-A}\left\vert x_2-x_1\right\vert<\varepsilon$. It follows that for integers $m,n$ with $n\ge m\ge N+1$,
\begin{align*}
\left\vert x_n-x_m\right\vert
&\le\left\vert x_{m+1}-x_m\right\vert+\left\vert x_{m+2}-x_{m+1}\right\vert+\cdots+\left\vert x_n-x_{n-1}\right\vert \\
&\le\left(A^{m-1}+A^m+\cdots+A^{n-2}\right)\left\vert x_2-x_1\right\vert \\
&\le \frac{A^{m-1}}{1-A}\left\vert x_2-x_1\right\vert \\
&\le \frac{A^{N}}{1-A}\left\vert x_2-x_1\right\vert \\
&<\varepsilon
\end{align*}
i.e., $\{x_n\}$ forms a Cauchy sequence in $\mathbb{R}$, and then $\lim_{n\to\infty}x_n=x$ for some $x\in \mathbb{R}$. Note that such $x$ is the desired fixed point of $f$, since
\begin{align*}
f(x)
=\lim_{n\to\infty}f(x_n)
=\lim_{n\to\infty}x_{n+1}
=x
\end{align*}

\smallskip
\noindent\textbf{Solution.}\quad
(d) The visualization is clear.



\bigskip\noindent\hrule\bigskip

The function $f$ defined by
\begin{align*}
f(x)=\frac{x^3+1}{3}
\end{align*}
has three fixed points, say $\alpha$, $\beta$, $\gamma$, where
\begin{align*}
-2<\alpha<-1,\quad
0<\beta<1,\quad
1<\gamma<2
\end{align*}
For arbitrarily chosen $x_1$, define $\{x_n\}$ by setting $x_{n+1}=f(x_n)$.
(a) If $x_1<\alpha$, prove that $x_n\to-\infty$ as $n\to\infty$.
(b) If $\alpha<x_1<\gamma$, prove that $x_n\to\beta$ as $n\to\infty$.
(c) If $\gamma<x_1$, prove that $x_n\to+\infty$ as $n\to\infty$.
Thus $\beta$ can be located by this method, but $\alpha$ and $\gamma$ cannot.

\smallskip
\noindent\textbf{Solution.}\quad

Define $g(x)=f(x)-x$, then $g(\alpha)=g(\beta)=g(\gamma)=0$ and
\begin{align*}
g(-2)=-\frac{1}{3},\quad g(-1)=1,\quad g(0)=\frac{1}{3},\quad g(1)=-\frac{1}{3},\quad g(2)=1
\end{align*}
We claim that 
(1) $g(x)<0$ for $x<\alpha$;
(2) $g(x)>0$ for $\alpha<x<\beta$;
(3) $g(x)<0$ for $\beta<x<\gamma$;
(4) $g(x)>0$ for $\gamma<x$.
\textit{Proof of the claim.} We only prove (1) because the others are in a similar way. Since $g$ is a polynomial of degree $3$ and $g$ has three zeros $\alpha,\beta,\gamma$, it is impossible that $g(x)=0$ whenever $x<\alpha$. Now, suppose there exists $x_0<\alpha$ such that $g(x_0)>0$, then $g(-2)<0<g(x_0)$ implies there exists $c$ between $-2$ and $x_0$, such that $g(c)=0$. Since $-2<\alpha$ and $x_0<\alpha$, we have $c<\alpha$, which is impossible to happen. $\Box$
To complete the proof, we consider the following five cases.
Case 1: $x_1<\alpha$. Suppose $x_n<\alpha$, then $x_{n+1}=\frac{x_n^3+1}{3}<\frac{\alpha^3+1}{3}=\alpha$ (since $x^3$ monotonically increases). So by mathematical induction, $x_n<\alpha$ for all $n$. It follows from (1) that $x_{n+1}=g(x_n)+x_n<x_n$ for each $n$, i.e., $\{x_n\}$ is monotonically decreasing. Suppose $\{x_n\}$ is bounded below, then $\lim_{n\to\infty}x_n=x'$ for some $x'\in \mathbb{R}$, and 
\begin{align*}
f(x')=\lim_{n\to\infty}f(x_n)=\lim_{n\to\infty}x_{n+1}=x'
\end{align*}
i.e., $x'$ is a fixed point of $f$. Since $g(x')=0$ and $x'\le x_1<\alpha$, a contradiction occurs (since $g$ has only three zeros). Hence $\lim_{n\to\infty}x_n=-\infty$.
Case 2: $\alpha<x_1<\beta$. Suppose $\alpha<x_n<\beta$, then $x_{n+1}=\frac{x_n^3+1}{3}>\frac{\alpha^3+1}{3}=\alpha$ and $x_{n+1}=\frac{x_n^3+1}{3}<\frac{\beta^3+1}{3}=\beta$ (since $x^3$ monotonically increases). So by mathematical induction, $\alpha<x_n<\beta$ for all $n$. It follows from (2) that $x_{n+1}=g(x_n)+x_n>x_n$ for each $n$, i.e., $\{x_n\}$ is monotonically increasing. Since $\{x_n\}$ is bounded above (by $\beta$), we have $\lim_{n\to\infty}x_n=x'$ for some $x'\in \mathbb{R}$, and 
\begin{align*}
f(x')=\lim_{n\to\infty}f(x_n)=\lim_{n\to\infty}x_{n+1}=x'
\end{align*}
i.e., $x'$ is a fixed point of $f$. Since $\alpha<x_1\le x'\le\beta$, $g(x')=0$, and $g(x)>0$ for all $x\in(\alpha,\beta)$, this means that $x'=\beta$.
Case 3: $x_1=\beta$. It is clear that $x_n=\beta$ for all $n$, implying $\lim_{n\to\infty}x_n=\beta$.
Case 4: $\beta<x_1<\gamma$. Similar to Case 2, the conclusion is $\lim_{n\to\infty}x_n=\beta$.
Case 5: $\gamma<x_1$. Similar to Case 1, the conclusion is $\lim_{n\to\infty}x_n=+\infty$.
In summary, part (a) follows by Case 1; part (b) follows by Case 2, 3, and 4; part (c) follows by Case 5.



\bigskip\noindent\hrule\bigskip

The process described in part (c) of Exercise 22 can of course also be applied to functions that map $(0,\infty)$ to $(0,\infty)$.
Fix some $\alpha>1$, and put
\begin{align*}
f(x)=\frac{1}{2}\left(x+\frac{\alpha}{x}\right),\quad g(x)=\frac{\alpha+x}{1+x}
\end{align*}
Both $f$ and $g$ have $\sqrt{\alpha}$ as their fixed point in $(0,\infty)$. Try to explain, on the basis of properties of $f$ and $g$, why the convergence in Exercise 16, Chap. 3, is so much more rapid than it is in Exercise 17. (Compare $f'$ and $g'$, draw the zig-zags suggested in Exercise 22.)

\smallskip
\noindent\textbf{Solution.}\quad

Let $r=\sqrt{\alpha}$.  For the first iteration,
\[
 f(x)-r=\frac{(x-r)^2}{2x}.
\]
Thus, once $x$ is near $r$, the new error is proportional to the square of the old error.  Equivalently, $f'(r)=0$, so the iteration has quadratic convergence.

For the second iteration,
\[
 g(x)-r=\frac{1-r}{1+x}(x-r),
 \qquad
 g'(r)=\frac{1-\alpha}{(1+r)^2}=\frac{1-r}{1+r}\ne0.
\]
Its error is therefore reduced only by an approximately fixed factor $|(1-r)/(1+r)|<1$ at each step: the convergence is linear.  The zig-zag diagram reflects the same distinction.  The graph of $f$ is tangent to the line $y=r$ at the fixed point, whereas the graph of $g$ crosses it with nonzero slope.



\bigskip\noindent\hrule\bigskip

Suppose $f$ is twice differentiable on $[a,b]$, $f(a)<0$, $f(b)>0$, $f'(x)\ge\delta>0$, and $0\le f''(x)\le M$ for all $x\in[a,b]$. Let $\xi$ be the unique point in $(a,b)$ at which $f(\xi)=0$.
Complete the details in the following outline of \textit{Newton's method} for computing $\xi$.
(a) Choose $x_1\in(\xi,b)$, and define $\{x_n\}$ by
\begin{align*}
x_{n+1}=x_n-\frac{f(x_n)}{f'(x_n)}
\end{align*}
Interpret this geometrically, in terms of a tangent to the graph of $f$.
(b) Prove that $x_{n+1}\le x_n$ (c.f. Rudin's book says $x_{n+1}<x_n$, but sometimes "$=$" may hold. For example, consider $f(x)=cx+d$ where $c>0$, then $x_n=\xi=-d/c$ for $n=2,3,\ldots$.) and that
\begin{align*}
\lim_{n\to\infty}x_n=\xi
\end{align*}
(c) Use Taylor's theorem to show that
\begin{align*}
x_{n+1}-\xi=\frac{f''(t_n)}{2f'(x_n)}\left(x_n-\xi\right)^2
\end{align*}
for some $t_n\in(\xi,x_n)$.
(d) If $A=M/2\delta$, deduce that
\begin{align*}
0\le x_{n+1}-\xi\le\frac{1}{A}\left[A\left(x_1-\xi\right)\right]^{2^n}
\end{align*}
(Compare with Exercise 16 and 18, Chap. 3.)
(e) Show that Newton's method amounts to finding a fixed point of the function $g$ defined by
\begin{align*}
g(x)=x-\frac{f(x)}{f'(x)}
\end{align*}
How does $g'(x)$ behave for $x$ near $\xi$?
(f) Put $f(x)=x^{1/3}$ on $(-\infty,\infty)$ and try Newton's method. What happens?

\smallskip
\noindent\textbf{Solution.}\quad

(a) Note that the tangent line of the graph of $f$ passing through $(x_n,f(x_n))$ is of the form
\begin{align*}
y-f(x_n)=f'(x_n)\left(x-x_n\right)
\end{align*}
Thus $x_{n+1}=x_n-\frac{f(x_n)}{f'(x_n)}$ is the intersection of the tangent line and the $x$-axis.

\smallskip
\noindent\textbf{Solution.}\quad
(b) We first show that $x_{n+1}\le x_n$. Since $f''\ge 0$, the function $f'$ is increasing. If $x_n=\xi$ for some $n$, then clearly $x_m=\xi$ for all $m\ge n$ and it is nothing to prove. If $x_n>\xi$, then by the mean value theorem there exists $c_n\in(\xi,x_n)$ such that
\begin{align*}
f(x_n)=f(x_n)-f(\xi)=f'(c_n)\left(x_n-\xi\right)\le f'(x_n)\left(x_n-\xi\right)
\end{align*}
Since $f'\ge\delta>0$, this reveals that $\frac{f(x_n)}{f'(x_n)}\le x_n-\xi$, and that
\begin{align*}
x_{n+1}=x_n-\frac{f(x_n)}{f'(x_n)}\ge x_n-\left(x_n-\xi\right)=\xi
\end{align*}
Note that $x_1>\xi$, so by induction, we conclude that $x_n\ge\xi$ for all $n$, and then $f(x_n)\ge 0$ for all $n$ (since $f'>0$ implies $f$ is increasing). It follows that $\frac{f(x_n)}{f'(x_n)}\ge 0$, or equivalently $x_{n+1}\le x_n$ for all $n$.
We next show that $\lim_{n\to\infty}x_n=\xi$. Since $\{x_n\}$ is monotonically decreasing with a lower bound $\xi$, $\lim_{n\to\infty}x_n=x$ for some $x\ge\xi$. But then $x=x-\frac{f(x)}{f'(x)}$, implying $f(x)=0$, so by the uniqueness, $x=\xi$.

\smallskip
\noindent\textbf{Solution.}\quad
(c) By Taylor's theorem, there exists $t_n\in(\xi,x_n)$ such that
\begin{align*}
f(\xi)=f(x_n)+f'(x_n)\left(\xi-x_n\right)+\frac{f''(t_n)}{2}\left(\xi-x_n\right)^2
\end{align*}
Since $f(\xi)=0$, we have
\begin{align*}
x_{n+1}-\xi
&=x_n-\xi-\frac{f(x_n)}{f'(x_n)} \\
&=x_n-\xi+\frac{1}{f'(x_n)}\left[f'(x_n)\left(\xi-x_n\right)+\frac{f''(t_n)}{2}\left(\xi-x_n\right)^2\right] \\
&=\frac{f''(t_n)}{2f'(x_n)}\left(x_n-\xi\right)^2
\end{align*}

\smallskip
\noindent\textbf{Solution.}\quad
(d) By part (c), we have
\begin{align*}
x_{n+1}-\xi
&\le A\left(x_n-\xi\right)^2 \\
&\le A\cdot A^2\left(x_{n-1}-\xi\right)^4 \\
&\,\,\vdots \\
&\le A\cdot A^2\cdots A^{2^{n-1}}\left(x_1-\xi\right)^{2^n} \\
&=\frac{1}{A}\left[A\left(x_1-\xi\right)\right]^{2^n}
\end{align*}

\smallskip
\noindent\textbf{Solution.}\quad
(e) Since $g(x)=x$ if and only if $f(x)=0$, the result then follows. Next, compute
\begin{align*}
g'(x)=1-\frac{\left[f'(x)\right]^2-f(x)f''(x)}{\left[f'(x)\right]^2}=\frac{f(x)f''(x)}{\left[f'(x)\right]^2}
\end{align*}
It follows that $g(x)$ tends to $0$ as $x$ near $\xi$.

\smallskip
\noindent\textbf{Solution.}\quad
(f) Given $x_n\in(-\infty,\infty)$ with $x_n\ne 0$, then $f'(x_n)=\frac{1}{3}\,x_n^{-\frac{2}{3}}$. So
\begin{align*}
x_{n+1}
=x_n-\dfrac{x_n^{\frac{1}{3}}}{\dfrac{1}{3}\,x_n^{-\frac{2}{3}}}
=-2x_n
\end{align*}
and we see that $\{x_n\}$ oscillates and diverges.



\bigskip\noindent\hrule\bigskip

Suppose $f$ is differentiable on $[a,b]$, $f(a)=0$, and there is a real number $A$ such that $\left\vert f'(x)\right\vert\le A\left\vert f(x)\right\vert$ on $[a,b]$. Prove that $f(x)=0$ for all $x\in[a,b]$. \textit{Hint:} Fix $x_0\in[a,b]$, let
\begin{align*}
M_0=\sup\left\vert f(x)\right\vert,\quad M_1=\sup\left\vert f'(x)\right\vert
\end{align*}
for $a\le x\le x_0$. For any such $x$,
\begin{align*}
\left\vert f(x)\right\vert\le M_1(x_0-a)\le A(x_0-a)M_0
\end{align*}
Hence $M_0=0$ if $A(x_0-a)<1$. That is, $f=0$ on $[a,b]$. Proceed.

\smallskip
\noindent\textbf{Solution.}\quad

Following the hint, choose some $x_0\in(a,b]$ such that $A(x_0-a)<1$. If $M_0=0$, then clearly $f=0$, and we can proceed. If $M_0>0$, note that
\begin{align*}
\left\vert f'(x)\right\vert
\le A\left\vert f(x)\right\vert
\le AM_0
\end{align*}
for $x\in[a,b]$, so $M_1\le AM_0$. Now, denote $\delta=1-A(x_0-a)$, then for $x\in[a,x_0]$,
\begin{align*}
\left\vert f(x)\right\vert
=\left\vert f'(c)\right\vert(x-a)
\le M_1(x_0-a)
\le A(x_0-a)M_0
=M_0-\delta M_0
\end{align*}
where $c\in(a,x)$ exists because of the mean value theorem. It follows that $M_0-\delta M_0$ is a new upper bound of $
\left\vert f(x)\right\vert$ for $x\in[a,x_0]$, contradicting the definition of $M_0$. So $M_0>0$ is impossible, and we conclude that $f=0$ on $[a,x_0]$. Next, choose $x_k=(k+1)x_0-ka$ for $k=0,1,2,\ldots,n-1$, where $n\ge 1$ satisfying $x_n=b$ and $x_n-x_{n-1}\le x_0-a$. By using the same argument, $f=0$ on $[x_{k-1},x_k]$, for $k=1,2,\ldots,n$. Hence $f=0$ on $[a,b]$.



\bigskip\noindent\hrule\bigskip

Let $\phi$ be a real function defined on a rectangle $R$ in the plane, given by $a\le x\le b$, $\alpha\le y\le \beta$. A \textit{solution} of the initial-value problem
\begin{align*}
y'=\phi(x,y),\quad y(a)=c\quad(\alpha\le c\le\beta)
\end{align*}
is, by definition, a differentiable function $f$ on $[a,b]$ such that $f(a)=c$, $\alpha\le f(x)\le\beta$, and
\begin{align*}
f'(x)=\phi(x,f(x))\quad(a\le x\le b)
\end{align*}
Prove that such a problem has at most one solution if there is a constant $A$ such that
\begin{align*}
\left\vert\phi(x,y_2)-\phi(x,y_1)\right\vert\le A\left\vert y_2-y_1\right\vert
\end{align*}
whenever $(x,y_1)\in R$ and $(x,y_2)\in R$.
\textit{Hint:} Apply Exercise 26 to the difference of two solutions. Note that this uniqueness theorem does not hold for the initial-value problem
\begin{align*}
y'=y^{1/2},\quad y(0)=0
\end{align*}
which has two solutions: $f(x)=0$ and $f(x)=x^2/4$. Find all other solutions.

\smallskip
\noindent\textbf{Solution.}\quad

(i) Let $f_1$ and $f_2$ be two solutions of the initial-value problem, define the difference function $f=f_2-f_1$ on $[a,b]$. It suffices to show that $f=0$ on $[a,b]$. Observe that $f(a)=f_2(a)-f_1(a)=c-c=0$, and that
\begin{align*}
\left\vert f'(x)\right\vert
&=\left\vert f_2'(x)-f_1'(x)\right\vert \\
&=\left\vert \phi(x,f_2(x))-\phi(x,f_1(x))\right\vert \\
&\le A\left\vert f_2(x)-f_1(x)\right\vert \\
&=A\left\vert f(x)\right\vert
\end{align*}
for all $x\in[a,b]$. Thus by Exercise 26, $f=0$ on $[a,b]$.
(ii) It is easy to check that $f(x)=0$ and $f(x)=\frac{x^2}{4}$ are solutions of the given initial-value problem. To find the others, observe that if $f$ is a \textit{nonzero} solution, then $y'=y^{\frac{1}{2}}$ implies $f'(x)=\left[f(x)\right]^{\frac{1}{2}}$. Differentiate it we get
\begin{align*}
f''(x)
=\frac{1}{2}\left[f(x)\right]^{-\frac{1}{2}}f'(x)
=\frac{1}{2}\left[f(x)\right]^{-\frac{1}{2}}\left[f(x)\right]^{\frac{1}{2}}
=\frac{1}{2}
\end{align*}
So $f'(x)=\frac{x}{2}+c$ for some constant $c$, and then $f(x)=\left(\frac{x}{2}+c\right)^2$. Since $y(0)=0$ implies $f(0)=0$, we get $c=0$. Hence $f(x)=\frac{x^2}{4}$.



\bigskip\noindent\hrule\bigskip

Formulate and prove an analogous uniqueness theorem for systems of differential equations of the form
\begin{align*}
y_j'=\phi_j(x,y_1,\ldots,y_k),\quad
y_j(a)=c_j,\quad
(j=1,\ldots,k)
\end{align*}
Note that this can be rewritten in the form
\begin{align*}
\mathbf{y}'=\Phi(x,\mathbf{y}),\quad\mathbf{y}(a)=\mathbf{c}
\end{align*}
where $\mathbf{y}=(y_1,\ldots,y_k)$ ranges over a $k$-cell, $\Phi$ is the mapping of a $(k+1)$-cell into the Euclidean $k$-space whose components are the functions $\phi_1,\ldots,\phi_k$, and $\mathbf{c}$ is the vector $(c_1,\ldots,c_k)$. Use Exercise 26, for vector-valued functions.

\smallskip
\noindent\textbf{Solution.}\quad

Statement: \textit{Let $\Phi$ be a vector-valued mapping of a $(k+1)$-cell $C=[a,b]\times I$ into $\mathbb{R}^k$. Suppose there exists a constant $A$ such that
\begin{align*}
\left\vert\Phi(x,\mathbf{y}_2)-\Phi(x,\mathbf{y}_1)\right\vert\le A\left\vert\mathbf{y}_2-\mathbf{y}_1\right\vert
\end{align*}
whenever $(x,\mathbf{y}_1)\in C$ and $(x,\mathbf{y}_2)\in C$. Then the initial-value problem
\begin{align*}
\mathbf{y}'=\Phi(x,\mathbf{y}),\quad \mathbf{y}(a)=\mathbf{c}\quad(\mathbf{c}\in I)
\end{align*}
has at most one solution.}
To show this, let $\mathbf{f}_1$ and $\mathbf{f}_2$ be two solution of the given initial-value problem, define $\mathbf{f}=\mathbf{f}_2-\mathbf{f}_1$ on $[a,b]$. It suffices to show that $\mathbf{f}=\mathbf{0}$ on $[a,b]$. Observe that $\mathbf{f}(a)=\mathbf{f}_2(a)-\mathbf{f}_1(a)=\mathbf{c}-\mathbf{c}=\mathbf{0}$, and that
\begin{align*}
\left\vert \mathbf{f}'(x)\right\vert
&=\left\vert \mathbf{f}_2'(x)-\mathbf{f}_1'(x)\right\vert \\
&=\left\vert \Phi(x,\mathbf{f}_2(x))-\Phi(x,\mathbf{f}_1(x))\right\vert \\
&\le A\left\vert \mathbf{f}_2(x)-\mathbf{f}_1(x)\right\vert \\
&=A\left\vert \mathbf{f}(x)\right\vert
\end{align*}
for all $x\in[a,b]$. Thus by Exercise 26 for vector-valued functions, $\mathbf{f}=\mathbf{0}$ on $[a,b]$.



\bigskip\noindent\hrule\bigskip

Specialize Exercise 28 by considering the system
\begin{align*}
y_j'=y_{j+1}\quad(j=1,\ldots,k-1)\\
y_k'=f(x)-\sum_{j=1}^kg_j(x)y_j
\end{align*}
where $f,g_1,\ldots,g_k$ are continuous real functions on $[a,b]$, and derive a uniqueness theorem for solutions of the equation
\begin{align*}
y^{(k)}+g_k(x)y^{(k-1)}+\cdots+g_2(x)y'+g_1(x)y=f(x)
\end{align*}
subject to initial conditions
\begin{align*}
y(a)=c_1,\quad
y'(a)=c_2,\quad
\ldots,\quad
y^{(k-1)}(a)=c_k
\end{align*}

\smallskip
\noindent\textbf{Solution.}\quad

Put $\Phi(x,y_1,y_2,\ldots,y_k)=\left(y_2,y_3,\ldots,y_k,f(x)-\sum_{j=1}^kg_j(x)y_j\right)$ and $\mathbf{c}=(c_1,c_2,\ldots,c_k)$, then the given system coincides with the initial-value problem
\begin{align*}
\mathbf{y}'=\Phi(x,\mathbf{y}),\quad\mathbf{y}(a)=\mathbf{c}
\end{align*}
Observe that for $\mathbf{y}_1=(y_{11},y_{21},\ldots,y_{k1})$ and $\mathbf{y}_2=(y_{12},y_{22},\ldots,y_{k2})$,
\begin{align*}
\left\vert\Phi(x,\mathbf{y}_2)-\Phi(x,\mathbf{y}_1)\right\vert^2
&=\sum_{i=2}^k(y_{i2}-y_{i1})^2+\left[\sum_{j=1}^kg_j(x)(y_{j2}-y_{j1})\right]^2
\end{align*}
Denote $M=\sup\{\left\vert g_j(x)\right\vert:x\in[a,b],\,1\le j\le k\}$, then
\begin{align*}
\left\vert\Phi(x,\mathbf{y}_2)-\Phi(x,\mathbf{y}_1)\right\vert^2
&\le\sum_{i=2}^k\left(y_{i2}-y_{i1}\right)^2+M^2\left[\sum_{j=1}^k\left(y_{j2}-y_{j1}\right)\right]^2 \\
&\le\sum_{i=2}^k\left(y_{i2}-y_{i1}\right)^2+kM^2\sum_{j=1}^k\left(y_{j2}-y_{j1}\right)^2 \\
&\le\left(1+kM^2\right)\sum_{j=1}^k\left(y_{j2}-y_{j1}\right)^2 \\
&=\left(1+kM^2\right)\left\vert\mathbf{y}_2-\mathbf{y}_1\right\vert^2
\end{align*}
Hence the uniqueness of the solution follows.


\end{document}
